%% ============================================================
%% Chapter 5 Solutions
%% ============================================================

\section{Chapter 5: The Adjoint in Linear Algebra}

%% ----------------------------------------------------------
\begin{solution}{5.1}{L2}

\solanswer

\textbf{The three-step adjoint strategy in plain words:}

\begin{enumerate}
\item Solve the forward problem to find the state. Store the solution.
\item Solve one backward (adjoint) problem for each response functional,
      using the matrix transposed from the forward problem.
\item For each parameter, compute the sensitivity as a dot product of
      the adjoint solution with the parameter's contribution vector.
      No additional solves are needed.
\end{enumerate}

\textbf{Why the adjoint equation $A^T\vec{v} = \vec{c}$ does not depend on $q_j$:}

The adjoint equation encodes ``which directions in state space affect
the response $J = \vec{c}^T\vec{u}$?'' This question is entirely
about $J$ (encoded in $\vec{c}$) and the forward operator (encoded in $A$).
It does not involve any particular parameter $q_j$, because the adjoint
does not ask what the state is — only which state directions matter for $J$.
The parameter-dependence enters only in the dot product step
($\partial J/\partial q_j = \vec{v}^T \partial\vec{b}/\partial q_j$),
which costs just one inner product per parameter once $\vec{v}$ is known.

\textbf{Why this makes the adjoint efficient:}
With $m$ parameters, the forward approach requires $m$ forward solves
(one FSE per parameter). The adjoint requires only 1 adjoint solve
(independent of $m$), then $m$ dot products. For large $m$, the
savings are a factor of $m$.

\solnote{The key conceptual point: the adjoint asks about $J$ and $A$,
not about specific parameters. Students often confuse the adjoint solve
with a parameter-dependent computation. Full credit requires explaining
why $A^T\vec{v} = \vec{c}$ has no $q_j$ on either side.}

\end{solution}

%% ----------------------------------------------------------
\begin{solution}{5.2}{L3}

\solanswer

System: $A\vec{u} = \vec{b}$ with $A = \bigl(\begin{smallmatrix}3&1\\1&2\end{smallmatrix}\bigr)$,
$\vec{b} = (q, 2q)^T$, $J = u_1 + u_2$ ($\vec{c} = (1,1)^T$).

\solpart{a} \textbf{Forward solution at $q = 1$.}
$A^{-1} = \frac{1}{5}\bigl(\begin{smallmatrix}2&-1\\-1&3\end{smallmatrix}\bigr)$,
so $\vec{u} = A^{-1}\vec{b} = \frac{1}{5}(2\cdot1 - 1\cdot2, -1\cdot1+3\cdot2)^T
= \frac{1}{5}(0,5)^T = (0, 1)^T$.
Hence $u_1 = 0$, $u_2 = 1$.

\solpart{b} \textbf{FSE for $\partial\vec{u}/\partial q$.}
Differentiate $A\vec{u} = \vec{b}$ with respect to $q$:
$A\,(\partial\vec{u}/\partial q) = \partial\vec{b}/\partial q = (1,2)^T$.
Solve: $\partial\vec{u}/\partial q = A^{-1}(1,2)^T
= \frac{1}{5}(2-2,-1+6)^T = \frac{1}{5}(0,5)^T = (0, 1)^T$.

\solpart{c} \textbf{Adjoint equation.}
$A^T\vec{v} = \vec{c}$. Since $A$ is symmetric, $A^T = A$, so
$A\vec{v} = (1,1)^T$:
$\vec{v} = A^{-1}(1,1)^T = \frac{1}{5}(2-1,-1+3)^T = \frac{1}{5}(1,2)^T$.

\solpart{d} \textbf{Sensitivity via adjoint formula.}
$\partial J/\partial q = \vec{v}^T\,(\partial\vec{b}/\partial q)
= \frac{1}{5}(1,2)\cdot(1,2)^T = \frac{1}{5}(1+4) = 1$.

\solpart{e} \textbf{Direct verification.}
$J(q) = u_1(q) + u_2(q)$. From (b), $\partial J/\partial q
= \vec{c}^T(\partial\vec{u}/\partial q) = (1,1)\cdot(0,1)^T = 1$. $\checkmark$

\solnote{The system has $\vec{b} = q(1,2)^T$, so $J = u_1+u_2$ is linear in $q$
and $\partial J/\partial q = 1$ exactly. Students should note that both
methods (FSE and adjoint) give the same answer — this is Unifying Idea~2
at work.}

\end{solution}

%% ----------------------------------------------------------
\begin{solution}{5.3}{L3}

\solanswer

The decision rule: use FSE when $m < r$; use adjoint when $r < m$.
FSE requires $m$ forward solves (plus 1 baseline); adjoint requires $r$
adjoint solves (plus 1 forward).

\solpart{a} $m = 5$ parameters, $r = 20$ functionals.
FSE: $m + 1 = 6$ solves. Adjoint: $r + 1 = 21$ solves.
\textbf{Use FSE} ($6 < 21$).

\solpart{b} $m = 100$ parameters, $r = 3$ functionals.
FSE: $101$ solves. Adjoint: $4$ solves.
\textbf{Use adjoint} ($4 \ll 101$). Speedup factor: $\approx 25\times$.

\solpart{c} $m = 8$ parameters, $r = 8$ functionals.
FSE: 9 solves. Adjoint: 9 solves. Both require the same cost.
\textbf{Either works}; choose based on implementation availability.
The crossover is at $m = r$.

\solpart{d} $m = 1{,}000$ parameters, $r = 1$ functional, $t = 0.1$ s/solve.
FSE: $1001 \times 0.1 = 100.1$ s. Adjoint: $2 \times 0.1 = 0.2$ s.
\textbf{Use adjoint.} Speedup: $100.1/0.2 \approx \mathbf{500\times}$.

\solnote{Part (d) illustrates the regime where the adjoint is not just
convenient but essential: 1000 parameters with FSE would take the
same time as with any other parameter-based approach, while the adjoint
reduces it to 2 solves regardless of $m$. This is the core motivation
for Chapters 5--6.}

\end{solution}

%% ----------------------------------------------------------
\begin{solution}{5.4}{L4}

\solanswer

The symbiotic equilibrium is $\bar{u}_1 = (1-a)/(1-ac) = 0.625$,
$\bar{u}_2 = (1-c)/(1-ac) = 0.75$ (from Chapter 4, Ex 4.4).
The equilibrium Jacobian $J_F$ and FSE system are as computed in Ex 4.4.

\solpart{a} \textbf{FSE approach: 3 solves for 6 sensitivities.}

The FSE gives $\partial\vec{u}/\partial p_j$ for each of the 3 parameters
($a$, $b$, $c$) by solving the same $2\times2$ linear system
$J_F\,(\partial\vec{u}/\partial p_j) = -\partial\vec{f}/\partial p_j$.
This is 3 solves, each giving $(\partial\bar{u}_1/\partial p_j,\,\partial\bar{u}_2/\partial p_j)$.

Numerical results at $a=0.5$, $b=2$, $c=0.4$
(from Ex 4.4 for $p = c$; deriving for $a$ and $b$ similarly):

\begin{center}
\begin{tabular}{lccc}
\toprule
& $\SI{a}{\cdot}$ & $\SI{b}{\cdot}$ & $\SI{c}{\cdot}$ \\
\midrule
$\bar{u}_1$ & $-0.250$ & $0$ & $+0.250$ \\
$\bar{u}_2$ & $+0.208$ & $0$ & $-0.417$ \\
\bottomrule
\end{tabular}
\end{center}

\textit{Derivation for $\partial/\partial a$:}
$-\partial\vec{f}/\partial a = (0, b\bar{u}_1\bar{u}_2/...)$... but $f_1 = \bar{u}_1(1-\bar{u}_1-a\bar{u}_2) = 0$, so $\partial f_1/\partial a = -\bar{u}_1\bar{u}_2 = -0.625(0.75) = -0.46875$.
$\partial f_2/\partial a = 0$ (only $f_1$ contains $a$).
Solve: $J_F(\partial\vec{u}/\partial a) = (0.46875, 0)^T$, giving
$\SI{a}{\bar{u}_1} = (a/\bar{u}_1)\partial\bar{u}_1/\partial a = (0.5/0.625)(-0.3125) = -0.250$.

\textit{For $b$:} Only $f_2 = b\bar{u}_2(1-\bar{u}_2-c\bar{u}_1) = 0$, so $\partial f_2/\partial b = \bar{u}_2(1-\bar{u}_2-c\bar{u}_1) = 0$ (at equilibrium). Thus $\partial\vec{u}/\partial b = \vec{0}$: $\SI{b}{\bar{u}_1} = \SI{b}{\bar{u}_2} = 0$.

\solpart{b} \textbf{Adjoint approach: 2 solves for 6 sensitivities.}

Two response functionals: $J_1 = \bar{u}_1$ ($\vec{c}_1 = (1,0)^T$)
and $J_2 = \bar{u}_2$ ($\vec{c}_2 = (0,1)^T$). Two adjoint solves:
$J_F^T\vec{v}_k = \vec{c}_k$ for $k = 1, 2$.
Then $\partial J_k/\partial p_j = \vec{v}_k^T(-\partial\vec{f}/\partial p_j)$ for each $j$.

\solpart{c} \textbf{Verification:} Both give the table above. $\checkmark$

\solpart{d} \textbf{Comparison.}
FSE uses 3 solves (one per parameter), adjoint uses 2 solves (one per \QOI).
Adjoint is better here because $r = 2 < m = 3$.
FSE would be preferred when $m < r$, e.g., for a model with 2 parameters
and 5 response functionals.

\solnote{The zero SIs for $b$ are a natural consequence of the equilibrium
structure: $b$ scales $f_2$ uniformly, so at equilibrium ($f_2 = 0$),
$\partial f_2/\partial b = 0$. This is a structural zero, not a numerical
artifact. Students who get nonzero SI for $b$ have an error in computing
$\partial\vec{f}/\partial b$.}

\end{solution}

%% ----------------------------------------------------------
\begin{solution}{5.5}{L4}

\solanswer

$R_0 = k\beta\nu / [(\nu+\mu)(\gamma+\mu)]$.

\solpart{a} \textbf{Computation graph.}
Intermediate nodes: $n_1 = k\beta$, $n_2 = n_1\nu$, $n_3 = \nu+\mu$,
$n_4 = \gamma+\mu$, $n_5 = n_3 n_4$, $R_0 = n_2/n_5$.

\solpart{b} \textbf{Forward mode for $\partial R_0/\partial\nu$.}
Propagate $\dot{\nu} = 1$ (seeding the derivative w.r.t.\ $\nu$):
$\dot{n}_1 = 0$, $\dot{n}_2 = n_1\cdot1 + \nu\cdot0 = k\beta$,
$\dot{n}_3 = 1$, $\dot{n}_4 = 0$, $\dot{n}_5 = \dot{n}_3 n_4 + n_3\dot{n}_4
= (\gamma+\mu)$.
$\partial R_0/\partial\nu = (\dot{n}_2 n_5 - n_2\dot{n}_5)/n_5^2
= (k\beta(\nu+\mu)(\gamma+\mu) - k\beta\nu(\gamma+\mu)) / n_5^2
= k\beta(\gamma+\mu)\mu / [(\nu+\mu)^2(\gamma+\mu)^2]
= k\beta\mu / [(\nu+\mu)^2(\gamma+\mu)]$.
Arithmetic operations: approximately 8 multiplications, 4 additions.

\solpart{c} \textbf{Reverse mode for all 5 partial derivatives.}
Seed $\bar{R}_0 = 1$. Backpropagate:
$\bar{n}_2 = \bar{R}_0/n_5 = 1/n_5$, $\bar{n}_5 = -\bar{R}_0 n_2/n_5^2 = -R_0/n_5$.
$\bar{n}_3 = \bar{n}_5 n_4$, $\bar{n}_4 = \bar{n}_5 n_3$.
$\bar{n}_1 = \bar{n}_2\nu$, $\bar{\nu} = \bar{n}_2 n_1 + \bar{n}_3$.
$\bar{k} = \bar{n}_1\beta$, $\bar{\beta} = \bar{n}_1 k$.
$\bar{\gamma} = \bar{n}_4$, $\bar{\mu} = \bar{n}_3 + \bar{n}_4$.
One backward pass gives all 5 derivatives simultaneously.

\solpart{d} \textbf{Normalized SI tornado plot.}
Using interpretable parameters $\tau_E = 1/\nu$, $\tau_I = 1/\gamma$, $L = 1/\mu$:
$\SI{k}{R_0} = 1$, $\SI{\beta}{R_0} = 1$, $\SI{\tau_E}{R_0} = -\tau_E/(\tau_E+L)$,
$\SI{\tau_I}{R_0} = \tau_I L/(\tau_I+L) \cdot (1/R_0\tau_I) \approx L/(L+\tau_I)$,
$\SI{L}{R_0} \approx 0$ (for $L \gg \tau_E, \tau_I$).
The tornado plot shows $k$ and $\beta$ dominant (SI $= 1$), with
$\tau_E$ and $\tau_I$ small for $L = 10{,}000$ days.

\solnote{The reverse-mode computation (part c) is exactly the adjoint
method applied to a scalar computation graph. This exercise bridges
the algebraic adjoint (Chapter 5) with automatic differentiation (Chapter 7).}

\end{solution}

%% ----------------------------------------------------------
\begin{solution}{5.6}{L5}

\solanswer

\textbf{Adjoint SA design for river water-quality model: $m = 40$, $r = 2$.}

\solpart{a} \textbf{Form of the adjoint equation.}
The river-quality model is a system of ODEs or PDEs. For a generic
ODE model $\dot{\vec{u}} = \vec{F}(\vec{u};\vec{p})$ with response
$J_k = \int_0^T g_k(\vec{u})\,dt + h_k(\vec{u}(T))$, the adjoint ODE for
response $k$ is:
\[
  -\dot{\vec{\lambda}}^{(k)} = J_F^T(\vec{u}(t))\,\vec{\lambda}^{(k)}
  - \nabla_{\vec{u}} g_k(\vec{u}(t)), \qquad
  \vec{\lambda}^{(k)}(T) = \nabla_{\vec{u}} h_k(\vec{u}(T)).
\]
Two adjoint equations (one per \QOI), both depending on the stored
forward solution $\vec{u}(t)$.

\solpart{b} \textbf{Number of adjoint solves: 2.}
One adjoint solve per response functional. Each solve runs backward
from $t = T$ to $t = 0$, giving $\vec{\lambda}^{(k)}(t)$.
Total cost: 1 forward solve (store) + 2 adjoint solves.
Compare with FSE: 1 + 40 forward solves.
The adjoint is $\approx 20\times$ cheaper.

\solpart{c} \textbf{Validation.}
Verify that the adjoint-computed $\partial J_k/\partial p_j$ matches
centered finite differences for all 40 parameters: use $h = 6\times10^{-6}|p_j^*|$
and check agreement to 5 significant figures. Also verify the adjoint
at a simpler test case with analytic sensitivity (e.g., a single-compartment
degradation model where $J$ and $\partial J/\partial p$ are known analytically).

\solpart{d} \textbf{Advising regulators.}
Produce two tornado plots (one per \QOI) showing all 40 parameters
sorted by $|\partial J_k/\partial p_j|$. Annotate each parameter with its
plain-language interpretation and the regulatory lever available:
``Reducing phosphorus loading at source $i$ by $10\%$ reduces downstream
concentration by approximately $2.3\%$.'' Identify the top 5 sources
for each \QOI and report whether the same parameters appear in both
tornado plots (co-optimization opportunity) or different ones
(trade-off between the two monitoring objectives).

\solpart{e} \textbf{Potential misleading results.}
If the river model has a threshold effect (e.g., dissolved oxygen below
a critical level triggers biological collapse), the adjoint — which
linearizes around the current state — may show small SI values for
parameters that would be critical if the threshold were crossed.
To detect: run the model at several nominal parameter sets spanning the
plausible range; if any set pushes the model near the threshold,
the adjoint at the current operating point is not representative.
Correction: report sensitivity results at multiple operating points
and flag which parameters bring the system close to critical thresholds.

\solnote{Full credit requires all five components with operational
specificity. The validation plan (part c) must include both a
numerical check (FD agreement) and an analytic check (simple test case).
Vague statements like ``validate the code'' are insufficient.}

\end{solution}

%% ----------------------------------------------------------
\begin{solution}{5.7}{L6}

\solanswer

\solpart{a} \textbf{Is the first sentence correct?}

Yes, with an important qualifier. Both FSE and adjoint compute
$\partial J/\partial p_j$ for all $j$, and the mathematical result
is identical: $\partial J/\partial p_j = \vec{v}^T\partial\vec{b}/\partial p_j$
(adjoint) $= \vec{c}^T(\partial\vec{u}/\partial p_j)$ (FSE), and
these two expressions are equal because $\vec{c}^T A^{-1} = \vec{v}^T$
(by definition of the adjoint). So the \emph{answer} is identical.
The adjoint is faster when $m \gg r$; if $m < r$, the FSE is actually
faster. ``Just faster'' is only correct in the typical regime $m \gg r$.

\solpart{b} \textbf{Is the second sentence correct?}

No, for at least two reasons.

\textit{Situation 1: The number of parameters grows.} If the model
is extended from $m = 10$ to $m = 200$ parameters (e.g., adding
spatially distributed source terms), the FSE cost grows as $200$
solves while the adjoint remains at $r + 1$ solves. At $m = 200$
and $r = 1$, switching from FSE to adjoint gives a $100\times$ speedup.

\textit{Situation 2: Multiple response functionals for a well-optimized
adjoint.} If a colleague needs to compute sensitivities for an
additional 5 \QOIs on the same model, the FSE cost is unchanged
(still $m + 1$ solves) while the adjoint adds only 5 more backward
solves. For models where solving the adjoint ODE is cheap compared
to evaluating $\partial\vec{f}/\partial p_j$ for all $j$, the adjoint
is the only viable option at large $m$.

\solpart{c} \textbf{When is FSE preferred?}

FSE is preferred when $m < r$: the number of parameters is less than
the number of response functionals. FSE costs $m + 1$ solves;
adjoint costs $r + 1$ solves.

\textit{Specific example where FSE is strictly better:}
$m = 3$ parameters (e.g., a simple SIR model with $k$, $\beta$, $\tau$),
$r = 50$ response functionals (sensitivity at 50 time points, or for
50 spatial locations). FSE: $4$ solves. Adjoint: $51$ solves.
FSE is $51/4 \approx 13\times$ faster.

\solnote{The key insight is that ``adjoint is always better'' is a
myth. The correct statement is: use the approach whose cost
($\min(m,r) + 1$ solves) is smaller. This is Table 5.1 in the text.
Students who say only ``adjoint is better when $m$ is large'' miss
the symmetric counterexample. Full credit requires a specific numerical
example in part (c).}

\end{solution}
